\scanpage{48}
% archive sheet number: 47
Cela, plus le fait que $t_G$ et $\Omega^\vee\otimes_{\mathbb Z_p}A$ sont de rang $1$ sur $A$ resp $A\otimes_{\unsure{\mathbb Z_p|\partial_p}}A$, et que $\rho$ est surjectif, prouve que $t_G$ \uline{est déduit de} $\Omega^\vee\otimes_{\mathbb Z_p}A$ par le changt de base $A\otimes_{\mathbb Z_p}A\mathrel{\overset{\rho_0}{\relbar}}A$, donc \deleted{\unsure{par suite}} (comme $\Omega^\vee\otimes_{\mathbb Z_p}A$ est déduit de $\Omega^\vee$ par $A\xrightarrow{\lambda\to\lambda\otimes 1}A\otimes_{\mathbb Z_p}A$) on obtient un isom can.
\[
t_G\simeq\Omega^\vee\qquad\text{i.e.}\qquad\omega_G\simeq\Omega
\]
et la suite exacte $(**)$ est la suite exacte canonique évidente, \inserted{d'où $(*)$ par dualité.} Comme cette filtration étend celle de $M_A\otimes_A k$, on a retrouvé ici le th. d'existence et d'unicité de Lubin-Tate, (\deleted{\unsure{en}}sous une forme moins forte d'ailleurs en ce qui concerne l'unicité, et sous l'hyp. complémentaire que $\exists$ puiss. div. sur $\pi A$.

Si on abandonne cette dernière hypothèse, et qu'on utilise la classification des déformations à isogénie près, on trouve de $=$ qu'on a une suite exacte canonique
\diagram{\begin{tikzcd}[ampersand replacement=\&,column sep=small,row sep=small]
0\arrow[r,dash]\&\omega_G\otimes_A K\arrow[r,dash]\&M_K\arrow[r]\arrow[d,equal]\&t_{G^*}\otimes_{\unsure{A|p}}K\arrow[r,dash]\&0\\
\&\&\Omega\otimes_{\mathbb Z_p}\deleted{\unsure{k}}K\&\&
\end{tikzcd}}
ou encore
\diagram{\begin{tikzcd}[ampersand replacement=\&,column sep=small,row sep=small]
0\arrow[r,dash]\&\omega_{G^*}\otimes_A K\arrow[r,dash]\&\overset{\vee}{M}_K\arrow[r]\arrow[d,equal]\&t_G\otimes_A K\arrow[r]\&0\\
\&\&\Omega^\vee\otimes_{\mathbb Z_p}K\&\&
\end{tikzcd}}
et on trouve encore que
\[
t_G\otimes_A K\simeq\overset{\vee}{M}_K\otimes_{\Omega^\vee\otimes_{\mathbb Z_p}K}(K\otimes K,\rho_0)\simeq\Omega^\vee\otimes_A K
\]
